\chap Experiments on learning

In this chapter, the dot matrix encoding of English language
characters will be considered in detail.
There are two important facets to consider in these experiments:
learning time, and error rate.
Learning time is the length of time in machine (or algorithm) time
units~$t$ which is required for the network to produce results which
are sufficiently accurate.
Error rate is calculated as the total error, or~$E$, which is the
sum of the square of the difference between the expected output and
actual output of the network.

Two variables tend to govern both the learning time and the error
rate.
These variables are the learning calculation parameter~$\varepsilon$ and
the number of \neus\ in the middle layer.
Both are analysed in detail below.

For the rest of this thesis, we will consider only~3 layer networks.
In \nn\ problems, the fourth layer is usually used to remove
rotational and translational problems with pattern recognition: as
the behaviour of this 3 layer network will be examined for behaviour
with rotated and other data, no fourth layer will be considered.
As well, the error threshold level was fixed at~$0.2$.
That is, a \neu\ which is supposed to have an output of~1 or~0 is
not flagged as being in error if its output is greater than or equal
to~$0.8$, or less than or equal to~$0.2$, respectively.
Some authors drive their networks to~$0.9$ and~$0.1$, but the extra
computational effort for a PC implementation would be considerable
(in many cases, most \neus\ were within~$0.1$ of their desired
value, with only a few being as far as~$0.2$ away).
The final invariant used here is the spread of the initial weights.
They were each randomly set to a value between~$+0.3$ and~$-0.3$.

The two learning calculation variables,~$\alpha$ and~$\varepsilon$,
can have a wide range of values.
Just {\sl how} wide appears to depend largely upon the problem.
For all experiments examined here, $\alpha$ was set to~$0.9$.
This doesn't pose a problem, because beyond values of a few tens of~$t$ the 
equation $\alpha e^{-t} \Delta w(t-1)$ drops to almost zero.
Hence~$\alpha$ is not really a significant factor in the
calculations beyond the very beginning, and most authors set~$\alpha = 0.9$.

So all learning equation parameter variations refer
to~$\varepsilon$.
The number of \neus\ in the second layer (referred to as layer~1, since
counting starts at~0,) was also widely varied
during experimentation.
We examine below the effects of these parameters' variations 
on learning time~$t$ and error rate~$E$.
A summary is included at the end of this chapter.

\sec Variation of Learning Time with Parameters

This experiment dealt with the length of time taken to fully train
the network with the error criteria set above.
However, there is a subtle difference to most authors approach to
this problem.

As initial weights varied randomly from~$+0.3$ to~$-0.3$, most experiments 
here were run~4 times in order to get more statistical accuracy.
Just under~80 simulations were performed, and the averages and standard
deviations calculated.
Both parameters of interest, $\varepsilon$ and \neus\ in the middle
layer, were varied.
$\varepsilon$ values were set at~$0.3$, $0.5$, $0.9$, and~$1.0$.
Configurations of layer~1 \neus\ that were used 
are~$15$, $20$, $35$, $50$, and~$70$.
Up to~4 tests were carried out with each one.

The results are tabulated and graphed below.
Where possible, (and it as in all but one case due to a lack of
data,) means and standard deviations (SD) are shown.
$$
\vbox{\tabskip=0pt \offinterlineskip
\def\tablerule{\noalign{\hrule}}
\halign{\strut#& \vrule#\tabskip=1em plus2em&
 \hfil#& \vrule width1.5pt#&
 \hfil#\hfil& \vrule#&
 \hfil#\hfil& \vrule#&
 \hfil#\hfil& \vrule#&
 \hfil#\hfil& \vrule#&
 \hfil#\hfil& \vrule#&
 \hfil#\hfil& \vrule#&
 \hfil#\hfil& \vrule#&
 \hfil#\hfil& \vrule#\tabskip=0pt\cr
\tablerule
&&{\ninerm neurons}&&\multispan{15}\hfil$\varepsilon$\hfil&\cr
\tablerule
&&{\ninerm in}&&\multispan3\hfil 0.3\hfil&&\multispan3\hfil 0.5\hfil&&
 \multispan3\hfil 0.9\hfil&&\multispan3\hfil 1.0\hfil&\cr
\tablerule
&&{\ninerm layer 1}&&{\ninerm Mean}&&{\ninerm SD}
 &&{\ninerm Mean}&&{\ninerm SD}
 &&{\ninerm Mean}&&{\ninerm SD}&&{\ninerm Mean}&&{\ninerm SD}&\cr
\noalign{\hrule height1.5pt}
&&15&&416&&33&&204&&25&&217&&50&&226&&--&\cr
\tablerule
&&20&&348&&24&&170&&3&&214&&62&&202&&32&\cr
\tablerule
&&35&&278&&12&&160&&8&&195&&23&&254&&60&\cr
\tablerule
&&50&&235&&10&&158&&2&&288&&46&&435&&172&\cr
\tablerule
&&70&&220&&15&&190&&14&&337&&46&&654&&210&\cr
\tablerule
}}
$$
\centerline{Figure 10. Table of mean learning times and standard deviations.}

\figure {11} {8} {Mean learning time versus number of \neus}

From the graph and table in figures~10 and~11, some trends become
clear: an increase in~$\varepsilon$ does not necessarily make the
network converge to a solution any faster.
In fact, as~$\varepsilon$ increases beyond about~$0.5$, the network
begins to become unstable during learning.

This phenomenon can be observed with the $-v$ option set in
fishNET.
With large~$\varepsilon$, some output \neus\ converge quickly
towards the desired result.
However, the others tend to move in the opposite direction to the
one desired.
After several tens of~$t$, the outputs suddenly begin to ``swap'',
and the \neus\ that had an output very close to the desired output
(some will be less than~$0.05$ away) begin to move away from the
desired result, and the outputs which were the opposite tend to
moves towards the correct output.
This see-saw effect can continue for several hundred~$t$, until
usually it will decrease in oscillation size, and all outputs will move 
slowly toward the desired result.

Before the final set of experiments quoted here were run, a previous
set of~$\varepsilon$ values was tried.
In this group were the values~$1.2$, $1.5$, and~$2.0$.
None of these values converged systematically in less than~$t =
5000$, which is the default limiting value.
This fact, and the standard deviations of the learning times examined
below, tends to indicate that~$\varepsilon = 1.0$ is near the boundary
of stable learning for this particular problem.

The oscillatory effect also seems to get worse with more \neus.
The network appears to have trouble learning for both a
large~$\varepsilon$ and a large number of intermediate layer \neus,
and the huge standard deviations in these cases point this out.
In fact, for the case where~$\varepsilon = 1.0$, and middle layer
\neus\ numbered~70, the learning time varied from 
between~$t=361$ and~$t=838$.
This is a very different result to the~$\varepsilon = 0.3$ 
or~$\varepsilon = 0.5$ standard deviations.
A graph of standard deviation versus \neus\ in layer~1 for
all~$\varepsilon$ values is shown in figure~12.
In general, for~$\varepsilon > 0.5$, the standard deviation of
learning time (and hence the instability during learning) increases
as a function of both~$\varepsilon$ and \neus\ in the middle layer.
\figure {12} {6} {Standard deviation versus number of \neus}

The effect of $\varepsilon$ and \neus\ in layer~1 on error rate will
be examined in the next section.

\sec Variation of Error Rate with Parameters

The experiment detailed here entailed taking a random sample of each
of the~$\varepsilon$ and middle layer values used in the previous
experiment and calculating the total error for a constant set of
input data.
Hence there are~20 values of total error.

The input data used is shown in the appendices.
It involves~4 sets of data for each of the~5 letters, including
noisy and distorted shapes.
None of the data used for testing had been used to train the network,
so the experiment was a good one to test the network's level of
generalisation.
Due to the large number of computations necessary, and the disk
space needed to store the networks, (training was done on the
Pyramid~9810, but all analysis, that is execution, was done on the
PC$_{8087}$,) only one of each category was tested.
This means that the shape of the graph of total error versus \neus\
should be viewed only for trends and may not itself be overly
statistically accurate.
The data appears below:
$$
\vbox{\tabskip=0pt \offinterlineskip
\def\tablerule{\noalign{\hrule}}
\halign{\strut#& \vrule#\tabskip=1em plus2em&
 \hfil#& \vrule width1.5pt#&
 \hfil#\hfil& \vrule#&
 \hfil#\hfil& \vrule#&
 \hfil#\hfil& \vrule#&
 \hfil#\hfil& \vrule#\tabskip=0pt\cr
\tablerule
&&{\ninerm neurons}&&\multispan{7}\hfil$\varepsilon$\hfil&\cr
\tablerule
&&{\ninerm (layer 1)}&&\hfil 0.3\hfil&&\hfil 0.5\hfil&&
 \hfil 0.9\hfil&&\hfil 1.0\hfil&\cr
\noalign{\hrule height1.5pt}
&&15&&6.15&&6.17&&5.61&&4.97&\cr
\tablerule
&&20&&6.03&&6.01&&5.57&&5.83&\cr
\tablerule
&&35&&6.17&&6.30&&5.46&&5.54&\cr
\tablerule
&&50&&6.11&&6.05&&5.80&&5.34&\cr
\tablerule
&&70&&6.25&&6.41&&5.81&&6.10&\cr
\tablerule
}}
$$
\centerline{Figure 13. Table of total errors.}

\figure{14} {6} {Total error versus \neus}

Figure~14 shows a graph of the data.
Interesting details can be drawn from it.
Firstly, the total error doesn't decrease as the number of \neus\ in
the intermediate layer grows larger.
Some authors suggest that the network fails to generalise when the
number of \neus\ grows too large.
The data sample here is too small to confirm or deny this, but there
seems to be a slight decrease, then increase in total error,
from left to right.

More interesting still is the effect of~$\varepsilon$ upon the
network.
Due to computational constraints, only a reasonably small data
sample was used, but it still appears that the error is a maximum
for some value of~$\varepsilon$ between about~$0.3$ and~$0.5$, and probably
gets smaller as~$\varepsilon$ gets smaller than~$0.5$ and also
smaller as~$\varepsilon$ gets larger than~$0.5$. 
(Note again that the
network would not reliably converge for values of~$\varepsilon$
above~$1.0$.)

Comparing the graphs in figures~11 and~14 yields even more important
results.
Networks that take longer to teach (converge) produce a smaller
error than those which converge quickly to a solution.
An explanation for this follows.
Faster converging solutions from the back-propagation algorithm are
more likely to settle in non-optimal minima in the solution space.
This is because the network will be ``coarser'' due to the obviously
larger values of~$\Delta w(t)$ that were added during
back-propagation calculations.
Because these values were larger, the network would probably not be
as finely tuned as one which took longer to converge from
smaller~$\Delta w(t)$s.
(Another way of thinking of this is to try to imagine that the
required solution is some symmetrical pattern of weights.
If the values of $\Delta w(t)$ are bigger, the network may satisfy
the error criteria but still not be particularly symmetric.)
Hence slower trained networks are more likely to accurately
generalise the problem, and produce fewer errors.

One final point to explain is why total error was used in this
experiment instead of actual error.
Even though none of the test data had been seen by the network, the
number of actual errors (that is incorrectly recognised letters)
{\sl was always the same for all cases above\/}, and was a very
small portion of the data.
(One can tell the selected pattern, even though the output for the
selected pattern's appropriate \neu\ is much less than the
desired~$0.8$, by looking for the most active \neu\ of the~5 output
\neus.)
In fact only the most distorted ``letters'' generated actual errors
(see the appendices), so total error was used to get some spread
between the various cases.

The next experiments examine the effect of casualties in the network
upon its performance.

\sec Summary

The results presented in this chapter highlight some features of
learning with this problem.
There can be large variations in the learning time for a particular
problem, depending upon the values of~$\varepsilon$ and the number
of middle layer \neus.
In general, learning time is shortest for values of~$\varepsilon$ of
between~$0.3$ and~$0.5$.
$\varepsilon$ values within this range cause networks to learn
faster as the number of \neus\ in the middle layer increases.
However, for larger values of~$\varepsilon$ (of between~$0.9$
and~$1.0$), the learning time {\sl increases\/} as the number of
\neus\ gets bigger.
A very large variance in learning times was observed for these
values.
This was especially true for~$\varepsilon = 1.0$, which to be on the
outer limits of convergence for large numbers of layer~1 \neus.

From looking at the error rate, we can make statements on the effect
of learning time on errors.
The error is slightly smaller for networks that take longer to
teach, and an explanation for this has been proposed.
There is not a dramatic difference in error rate between the
different~$\varepsilon$ and \neu\ values (about~20\% variation
at most).

